\section*{Chapter Summary}
Big data analytics frameworks, such as Spark and Giraph, need to process and cache massive datasets that do not always fit on the managed heap.
Therefore, frameworks temporarily move long-lived objects outside the heap (off-heap) on a fast storage device.
However, this practice results in (1) high serialization/deserialization (\teraheapSD) cost and (2) high memory pressure when off-heap objects are moved back for processing.

In this article, we propose \teraheapName{}, a system that eliminates \teraheapSD overhead and expensive \teraheapGC scans for a large portion of objects in analytics frameworks.
\teraheapName{} relies on three concepts: (1) It eliminates \teraheapSD by extending the managed runtime (JVM) to use a second high-capacity heap (H2) over a fast storage device.
(2) It offers a simple hint-based interface, allowing analytics frameworks to leverage object knowledge to populate H2.
(3) It reduces \teraheapGC cost by fencing the collector from scanning H2 objects while maintaining the illusion of a single managed heap, ensuring memory safety.

We implement \teraheapName{} in OpenJDK8 and OpenJDK17 and evaluate it with fifteen widely used applications in two real-world big data frameworks, Spark and Giraph.
We find that for the same DRAM size, \teraheapName{} improves performance by up to 73\% and 28\% compared to native Spark and Giraph.
Also, it can still provide better performance by consuming up to $4.6\times$ and $1.2\times$ less DRAM than native Spark and Giraph, respectively.
\teraheapName{} can also be used for in-memory frameworks and applying it to the Neo4j Graph Data Science library improves its performance by up to 26\%.
Finally, it outperforms Panthera, a state-of-the-art garbage collector for hybrid DRAM-NVM memories, by up to 69\%.
